\documentclass[12pt,a4paper,]{ctexart}
\usepackage[UTF8]{ctex}
\usepackage{amsmath,amssymb}
\usepackage{geometry}
\usepackage{xcolor}
\usepackage{listings}
\geometry{margin=2.5cm}
\pagestyle{empty}
\title{Narcissistic number(PPDI)}
\author{}
\date{}
\lstset{
	language=C++,
	basicstyle=\ttfamily\small,
	numbers=left,
	frame=single,
	tabsize=4,
	breaklines=true,
	showstringspaces=false,
	keywordstyle=\bfseries\color{red},
	directivestyle=\color{green},
	stringstyle=\color{blue},
	numberstyle=\tiny\color{gray}
}
\begin{document}
	\maketitle
	\begin{center}
		Narcissistic number (自幂数) has another name called Armstrong number (阿姆斯壯数)
	\end{center}
	\begin{flushleft}
		\section{1-digit Narcissistic number[独身数]}
		All number 0-9 is Narcissistic number.
		\section{2‑digit Narcissistic numbers}
		2‑digit Narcissistic numbers does not exist.
		\section{3-digit Narcissistic number(Daffodil number)[水仙花数]}
		This refers to a 3-digit number where the sum of the cube of its digits equals the number itself.\\
		Example: 153\\
		$1^3+5^3+3^3 = 153$\\
		So \underline{153} is a Daffodil number.\\
		There have 4 daffodil number are:
		\begin{enumerate}
			\item 153
			\item 370
			\item 371
			\item 407
		\end{enumerate}
		\subsection{Find daffodil number with programming}
		\subsubsection{In C++}
		\begin{lstlisting}
			#include<iostream>
			#include<cstdio>
			int main()
			{
				int a, b, c;
				for (int num = 100; num <= 999; num++) {
					a = num / 100;
					b = (num / 10) % 10;
					c = num % 10;
					if (a * a * a + b * b * b + c * c * c == num)
					printf("%d\n", num);
				}
			}
		\end{lstlisting}
		\subsubsection{In Python}
		\begin{lstlisting}
			for num in range(100,1000):
			a = num // 100
			b = (num // 10) % 10
			c = num % 10
			if a**3+b**3+c**3==num:
			print (num)
		\end{lstlisting}
		\section {4-digit Narcissistic number(Four-leaf-Rose Number)[四叶玫瑰数]}
		This refers to a 4-digit number where the sum of the fourth powers of its digits equals the number itself.\\
		Example: 1634\\
		$1^4+6^4+3^4+4^4 = 1634$\\
		So \underline{1634} is a Four-leaf-rose Number.\\
		There have 3 rose number are:
		\begin{enumerate}
			\item 1634
			\item 8208
			\item 9474
		\end{enumerate}
		\subsection {Find Four-leaf-rose Number with programming.}
		\subsubsection{In C++}
		\begin{lstlisting}
			#include<iostream>
			#include<cstdio>
			using namespace std;
			
			int main(){
				int a,b,c,d;
				for (int num=1000;num<=9999;num++){
					a = num / 1000;
					b = (num / 100) % 10;
					c = (num / 10) % 10;
					d = num % 10;
					if (a*a*a*a+b*b*b*b+c*c*c*c+d*d*d*d==num)
					printf("%d\n",num);
				}
				cin.git();
				return 0;
			}
		\end{lstlisting}
		\subsubsection{In Python}
		\begin{lstlisting}
			for num in range(1000,10000):
			a = num // 1000
			b = (num // 100) % 10
			c = (num // 10) % 10
			d = num % 10
			if a**4+b**4+c**4+d**4==num:
			print (num)
		\end{lstlisting}
		\section{5‑digit Narcissistic number(Five‑point‑star number)[五角星数]}
		This refers to a 5-digit number where the sum of the fiveth powers of its digits equals the number itself.\\
		There has 3 five-point-star number are:
		\begin{enumerate}
			\item 54748
			\item 92727
			\item 93084
		\end{enumerate}
		\subsection{Find five-point-star number with programing}
		\begin{lstlisting}
			#include<iostream>
			#include<cstdio>
			using namespace std;
			int main(){
				int a, b, c, d, e;
				for (int num = 10000; num <= 99999; num++) {
					a = num / 10000;
					b = (num / 1000) % 10;
					c = (num / 100) % 10;
					d = (num / 10) % 10;
					e = num % 10;
					if (a * a * a * a + b * b * b * b + c * c * c * c + d * d * d * d == num)
					printf("%d\n", num);
				}
				cin.get()
				return 0;
			}
		\end{lstlisting}
		\begin{lstlisting}
			for num in range(10000,100000):
			a = num // 10000
			b = (num // 1000) % 10
			c = (num // 100) % 10
			d = (num // 10) % 10
			e = num % 10
			if a**5+b**5+c**5+d**5+e**5==num:
			print (num)
		\end{lstlisting}
	\end{flushleft}
\end{document}
	

